
\subsection{$\mathcal M_{\mathrm{weak}}(C)$ 及其同伦不变性}
\label{subsec:4-3-2}

令 $C$ 为含幺滤过 $A_\infty$ 代数。回忆第 3.6.3 小节定义的集合
$\widehat{\mathcal M}_{\mathrm{weak}}(C)$。也就是说，
$b\in C^1=C[1]^0$ 属于 $\widehat{\mathcal M}_{\mathrm{weak}}(C)$，
当且仅当对某个 $c\in\Lambda^{+(0)}_{0,\mathrm{nov}}$ 有
\[
 \mathfrak m(e^b)=ce\mathbf e.
\]
这样的 $b$ 称为弱定界余链；见定义 3.6.29。还要回忆，含幺滤过
$A_\infty$ 同态 $\mathfrak f\colon C_1\to C_2$ 诱导映射
\[
 \mathfrak f_*\colon
 \widehat{\mathcal M}_{\mathrm{weak}}(C_1)
 \longrightarrow
 \widehat{\mathcal M}_{\mathrm{weak}}(C_2),
 \qquad
 \mathfrak f_*(b)=\mathfrak f(e^b).
\]
见引理 3.6.36。

令 $\mathfrak C$ 为 $[0,1]\times C$ 的一个含幺模型。

\begin{definition}{4.3.19}\label{def:4-3-19}
令 $b_0,b_1\in\widehat{\mathcal M}_{\mathrm{weak}}(C)$。若存在
$\widetilde b\in\widehat{\mathcal M}_{\mathrm{weak}}(\mathfrak C)$，使得
\[
 \mathrm{Eval}_{s=0*}(\widetilde b)=b_0,
 \qquad
 \mathrm{Eval}_{s=1*}(\widetilde b)=b_1,
\]
就称 $b_0$ 与 $b_1$ 规范等价，并记作 $b_0\sim b_1$。
这就是弱定界余链的规范等价\liftIndex{I155}。
\end{definition}

\begin{lemma}{4.3.20}\label{lem:4-3-20}
关系 $\sim$ 不依赖于 $[0,1]\times C$ 的模型之选择，并且是等价关系。
\end{lemma}

\begin{proof}
证明与引理 4.3.3、4.3.4 的证明相同。
\end{proof}

\begin{definition}{4.3.21}\label{def:4-3-21}
令 $(C,\mathfrak m)$ 为含幺滤过 $A_\infty$ 代数。定义
$\mathcal M_{\mathrm{weak}}(C,\mathfrak m)$ 为
$\widehat{\mathcal M}_{\mathrm{weak}}(C,\mathfrak m)$ 的所有
$\sim$-等价类组成的集合。
\end{definition}

下面的定理可以用与定理 4.3.13 相同的方法证明。

\begin{theorem}{4.3.22}\label{thm:4-3-22}
对应
$(C,\mathfrak m)\longmapsto\mathcal M_{\mathrm{weak}}(C,\mathfrak m)$
定义了协变函子
{\small
\[
 \begin{gathered}
 \mathfrak{HA}_{\infty u}:=
 [\text{含幺滤过 }A_\infty\text{ 代数；}\\[-0.2em]
 \text{含幺滤过 }A_\infty\text{ 同态的 u-同伦类}]
 \longrightarrow[\mathrm{set}].
 \end{gathered}
\]
}
若
$\mathfrak f\colon(C_1,\mathfrak m^1)\to(C_2,\mathfrak m^2)$
是含幺同伦等价，则它诱导双射
\[
 \mathfrak f_*\colon
 \mathcal M_{\mathrm{weak}}(C_1,\mathfrak m^1)
 \longrightarrow
 \mathcal M_{\mathrm{weak}}(C_2,\mathfrak m^2).
\]
\end{theorem}

第 3.6.3 小节通过
\[
 \mathfrak m(e^b)=\mathfrak{PO}(b)e\mathbf e
\]
引入了势函数
$\mathfrak{PO}\colon\widehat{\mathcal M}_{\mathrm{weak}}(C)
\to\Lambda^{+(0)}_{0,\mathrm{nov}}$（定义 3.6.33）。现在可以证明：

\begin{lemma}{4.3.23}\label{lem:4-3-23}
$\mathfrak{PO}(b)$ 只依赖于
$b\in\widehat{\mathcal M}_{\mathrm{weak}}(C)$ 的规范等价类。也就是说，
$\mathfrak{PO}$ 诱导映射
\[
 \mathfrak{PO}\colon
 \mathcal M_{\mathrm{weak}}(C)
 \longrightarrow\Lambda^{+(0)}_{0,\mathrm{nov}}.
\]
此外，若 $C_i$ 是含幺滤过 $A_\infty$ 代数，且
$\mathfrak f\colon C_1\to C_2$ 是含幺滤过 $A_\infty$ 同态，则对
$b\in\mathcal M_{\mathrm{weak}}(C_1)$ 有
\[
 \mathfrak{PO}_2\circ\mathfrak f_*(b)=\mathfrak{PO}_1(b),
\]
其中 $\mathfrak{PO}_i$ 是 $C_i$ 的势函数。
\liftIndex{I156}
\end{lemma}

\begin{proof}
令
$b_0\sim b_1$ 于 $\widehat{\mathcal M}_{\mathrm{weak}}(C)$，并令
$\widetilde b$ 为定义 4.3.19 所述从 $b_0$ 到 $b_1$ 的同伦
\translatornote{原书此处误写“定义 4.3.1”；弱定界余链的规范等价及这里
使用的含幺模型实际定义于定义 4.3.19。}。由引理 3.6.36 (3)，
\[
 \mathfrak{PO}(b_0)
 =\mathfrak{PO}(\mathrm{Eval}_{s=0*}\widetilde b)
 =\mathfrak{PO}(\widetilde b)
 =\mathfrak{PO}(\mathrm{Eval}_{s=1*}\widetilde b)
 =\mathfrak{PO}(b_1).
\]
这证明了第一个断言。第二个断言也由引理 3.6.36 (3) 得到。
\end{proof}

当 $C$ 是同伦含幺滤过 $A_\infty$ 代数时，定义
\begin{equation}
 \mathcal M_{\mathrm{weak}}(C)
 =\mathcal M_{\mathrm{weak}}(C^+).
\tag{4.3.24}\label{eq:4-3-24}
\end{equation}
这个集合在 $C$ 的 hu-同伦型下不变。势函数在这种情形中也用同样的方法
定义；它是定义在
$\mathcal M_{\mathrm{weak}}(C)=\mathcal M_{\mathrm{weak}}(C^+)$ 上的
函数。还可以用类似方式定义范畴 $\mathfrak{HA}_{\infty hu}$。

现在考虑滤过 $A_\infty$ 双模的情形。令 $C_0,C_1$ 为含幺滤过
$A_\infty$ 代数，令 $D$ 为左 $C_1$、右 $C_0$ 的含幺滤过
$A_\infty$ 双模（定义 3.7.5）。对
$b_0\in\widehat{\mathcal M}_{\mathrm{weak}}(C_0)$、
$b_1\in\widehat{\mathcal M}_{\mathrm{weak}}(C_1)$，定义--引理
3.7.13 定义了
$\delta_{b_1,b_0}\colon D\to D$，即
\[
 \delta_{b_1,b_0}(x)=\mathfrak n(e^{b_1}xe^{b_0}).
\]
以 $\mathfrak{PO}_i$（$i=0,1$）表示 $C_i$ 的势函数。命题 3.7.17
表明
\[
 \delta_{b_1,b_0}\circ\delta_{b_1,b_0}(x)
 =\bigl(-\mathfrak{PO}_1(b_1)+\mathfrak{PO}_0(b_0)\bigr)ex.
\]
因此，当 $\mathfrak{PO}_0(b_0)=\mathfrak{PO}_1(b_1)$ 时，可以定义
$(D,\delta_{b_1,b_0})$ 的上同调。

\begin{proposition}{4.3.25}\label{prop:4-3-25}
在上述情形下，令
\[
 b_0,b'_0\in\widehat{\mathcal M}_{\mathrm{weak}}(C_0),
 \qquad
 b_1,b'_1\in\widehat{\mathcal M}_{\mathrm{weak}}(C_1).
\]
假设
$\mathfrak{PO}_0(b_0)=\mathfrak{PO}_1(b_1)$
\translatornote{原书在这里省略势函数的下标，写成
$\mathfrak{PO}(b_0)=\mathfrak{PO}(b_1)$；但 $b_0,b_1$ 分属
$C_0,C_1$，本段此前已将相应势函数记为 $\mathfrak{PO}_0,\mathfrak{PO}_1$。}。
若 $b_0\sim b'_0$ 且 $b_1\sim b'_1$，则有典范同构
\[
 H(D,\delta_{b_1,b_0})
 \cong
 H(D,\delta_{b'_1,b'_0}).
\]
同伦含幺版本也成立。
\end{proposition}

证明将在第 5.2.4 小节给出（推论 5.2.40）。

\subsection{$\mathcal M_{\mathrm{weak,def}}(L)$ 及其同伦不变性}
\label{subsec:4-3-3}

第 3.8 节利用环境辛流形 $M$ 上的闭上链 $\mathfrak b$ 引入了无穷小形变
（bulk 形变）。回忆
\[
 \begin{aligned}
 \widehat{\mathcal M}_{\mathrm{weak,def}}(L)
 ={}&\bigl\{(\mathfrak b,b)\in
 C^2(M;\Lambda^+_{0,\mathrm{nov}})\\[-0.2em]
 &\quad{}\times C^1(L;\Lambda^+_{0,\mathrm{nov}})
 \ \bigm|\ 
 \delta_M\mathfrak b=0,\quad
 \mathfrak m_0^{+\mathfrak b,b}=ce\mathbf e^+
 \bigr\}.
 \end{aligned}
\]
见定义 3.8.40。

选定代表
$H^2(M;\Lambda^+_{0,\mathrm{nov}})$ 的子空间
\[
 \mathcal H^2(M;\Lambda^+_{0,\mathrm{nov}})
 \subset C(M;\Lambda^+_{0,\mathrm{nov}}).
\]

\begin{definition}{4.3.26}\label{def:4-3-26}
令
$(\mathfrak b,b_0),(\mathfrak b,b_1)
\in\widehat{\mathcal M}_{\mathrm{weak,def}}(L)$，其中
$\mathfrak b\in\mathcal H^2(M;\Lambda^+_{0,\mathrm{nov}})$。
若 $b_0$ 与 $b_1$ 在含幺滤过 $A_\infty$ 代数
\[
 \bigl(C(L;\Lambda_{0,\mathrm{nov}}),\mathfrak m^{\mathfrak b}\bigr)
\]
中依定义 4.3.19 规范等价，就称
$(\mathfrak b,b_0)$ 与 $(\mathfrak b,b_1)$ 规范等价，并记作
$(\mathfrak b,b_0)\sim(\mathfrak b,b_1)$\liftIndex{I157}。
\translatornote{原书在本定义及紧随其后的说明中把这个边界代数误写成
$(C(M;\Lambda^+_{0,\mathrm{nov}}),\mathfrak m^{\mathfrak b})$；但
$b_0,b_1$ 是 $L$ 上的边界余链，定义 3.8.40 给出的相应代数是
$C(L;\Lambda_{0,\mathrm{nov}})$。译文据此统一改正。}
于是 $\sim$ 显然定义了等价关系；再定义
$\mathcal M_{\mathrm{weak,def}}(L)$ 为所有这样的
$(\mathfrak b,b)$ 的规范等价类组成的集合。
\end{definition}

第 4.6.4 小节将证明
$\mathcal M_{\mathrm{weak,def}}(L)$ 不依赖于子空间
$\mathcal H^2(M;\Lambda^+_{0,\mathrm{nov}})$ 的选择。那里证明的定理
4.6.47 还蕴含：滤过 $A_\infty$ 代数
\[
 \bigl(C(L;\Lambda_{0,\mathrm{nov}}),\mathfrak m^{\mathfrak b}\bigr)
\]
在严格同伦等价的意义下只依赖于 $\mathfrak b$ 的上同调类
\translatornote{原书此处写“同调类”（homology class），并在下一句写
“同调”（homologous）；但 $\mathfrak b$ 是 $H^2$ 中的闭上链，式
(4.3.27) 也把它送到 $[\mathfrak b]\in H^2$，故均应为上同调概念。}。
进而，定义 4.6.50 将把定义 4.3.26 推广到
$(\mathfrak b_0,b_0)$、$(\mathfrak b_1,b_1)$ 这一情形，其中
$\mathfrak b_0$ 与 $\mathfrak b_1$ 代表同一个上同调类，但不必相等。

由定义，有良定义映射
\begin{equation}
 \pi_{\mathrm{amb}}\colon
 \mathcal M_{\mathrm{weak,def}}(L)
 \ni(\mathfrak b,b)
 \longmapsto[\mathfrak b]
 \in H^2(M;\Lambda^+_{0,\mathrm{nov}}),
\tag{4.3.27}\label{eq:4-3-27}
\end{equation}
这就是第 1 章定理 B 中提到的映射。现在有：

\begin{lemma}{4.3.28}\label{lem:4-3-28}
定义 3.8.40 中的势函数 $\mathfrak{PO}(\mathfrak b,b)$ 只依赖于定界余链
\[
 b\in\widehat{\mathcal M}_{\mathrm{weak}}
 \bigl(C(L;\Lambda_{0,\mathrm{nov}}),\mathfrak m^{\mathfrak b}\bigr)
\]
的规范等价类
\translatornote{原书此处把这一集合写成
$\widehat{\mathcal M}_{\mathrm{weak,def}}
(C(M;\Lambda^+_{0,\mathrm{nov}}),\mathfrak m^{\mathfrak b})$。
按定义 3.8.40，固定 $\mathfrak b$ 后，$b$ 属于 $L$ 上形变代数的
$\widehat{\mathcal M}_{\mathrm{weak}}$；译文据此同时改正底空间、系数环
和模空间记号。}。也就是说，$\mathfrak{PO}$ 诱导映射
\[
 \mathfrak{PO}\colon
 \mathcal M_{\mathrm{weak,def}}(L)
 \longrightarrow\Lambda^{+(0)}_{0,\mathrm{nov}}.
\]
\liftIndex{I156}
\end{lemma}

\begin{proof}
把引理 4.3.23 用于滤过 $A_\infty$ 代数
$\bigl(C(L;\Lambda_{0,\mathrm{nov}}),\mathfrak m^{\mathfrak b}\bigr)$
即可。
\end{proof}

第 4.6.4 小节还将证明 $\mathfrak{PO}$ 不依赖于
$\mathcal H^2(M;\Lambda^+_{0,\mathrm{nov}})$ 的选择（见引理
4.6.51）。关于势函数在辛微分同胚
$\psi\colon M\to M'$ 下的相应性质，见推论 4.6.53。

现在考虑第 3.8.7 小节所述的滤过 $A_\infty$ 双模，并沿用该小节的
记号。下面的结论是命题 4.3.18、4.3.25 的类似物，其证明也将在第
5.2.4 小节给出（推论 5.2.40）。

\begin{proposition}{4.3.29}\label{prop:4-3-29}
在第 3.8.7 小节的情形下，令
\[
 (\mathfrak b,b_0),(\mathfrak b,b'_0)
 \in\widehat{\mathcal M}_{\mathrm{weak,def}}(L^{(0)}),
 \qquad
 (\mathfrak b,b_1),(\mathfrak b,b'_1)
 \in\widehat{\mathcal M}_{\mathrm{weak,def}}(L^{(1)}).
\]
假设
\[
 \mathfrak{PO}(\mathfrak b,b_0)
 =\mathfrak{PO}(\mathfrak b,b_1).
\]
若
$(\mathfrak b,b_0)\sim(\mathfrak b,b'_0)$ 且
$(\mathfrak b,b_1)\sim(\mathfrak b,b'_1)$，则两个弗洛尔上同调之间有
典范同构
\[
 \begin{aligned}
 &HF\bigl(
 (L^{(1)},(\mathfrak b,b_1)),
 (L^{(0)},(\mathfrak b,b_0));
 \Lambda_{0,\mathrm{nov}}
 \bigr)\\
 &\qquad\cong
 HF\bigl(
 (L^{(1)},(\mathfrak b,b'_1)),
 (L^{(0)},(\mathfrak b,b'_0));
 \Lambda_{0,\mathrm{nov}}
 \bigr).
 \end{aligned}
\]
\end{proposition}

这个命题蕴含弗洛尔上同调是良定义的，并且由如下规范等价类之对子
\[
 (\mathfrak b,b_0),(\mathfrak b,b_1)
 \in\mathcal M_{\mathrm{weak,def}}(L)
\]
参数化，其中
\[
 \pi_{\mathrm{amb}}(\mathfrak b,b_0)
 =\pi_{\mathrm{amb}}(\mathfrak b,b_1),
 \qquad
 \mathfrak{PO}(\mathfrak b,b_0)
 =\mathfrak{PO}(\mathfrak b,b_1).
\]
这正是第 1 章定理 B 的 (B.1)。

从现在起，只要不会引起混淆，就用同一个符号
$(\mathfrak b,b)\in\mathcal M_{\mathrm{weak,def}}(L)$ 表示
$(\mathfrak b,b)\in\widehat{\mathcal M}_{\mathrm{weak,def}}(L)$ 的规范等价类。
